\section{Analytical Cost Model}
\label{sec:cost_models}
The empirical benchmarking uses static configurations swept explicitly to gather data on cost behavior. While theoretical "dynamic batching cost models" (which algorithmically predict optimal batch sizes based on fluctuating L1 base fees) are discussed in early project proposals, they are analytical constructs and are explicitly not implemented in the active runtime codebase.

Following the methodology of Chaliasos et al. \cite{chaliasos2024analyzing}, the cost model utilized for analyzing results splits into fixed batch costs (L1 settlement base fees, proof verification) and marginal transaction costs (DA bytes, witness generation, execution). 

\begin{table}[H]
\centering
\caption{ZK-Rollup Cost Model \cite{chaliasos2024analyzing, saif2024dataavailability}}
\begin{tabular}{|l|l|l|l|}
\hline
\textbf{Cost Category} & \textbf{Examples} & \textbf{Fixed/Marginal} & \textbf{Source} \\ \hline
L1 Settlement & Commit batch, base fee & Fixed per batch & Chaliasos \cite{chaliasos2024analyzing} \\ \hline
Proof Verification & Groth16 pairing check & Fixed per batch & Chaliasos \cite{chaliasos2024analyzing} \\ \hline
Data Availability (DA) & calldata, blob bytes & Marginal / payload-dependent & Saif \cite{saif2024dataavailability} \\ \hline
Proving Time & witness + proof generation & Marginal / cycles & Chaliasos \cite{chaliasos2024analyzing} \\ \hline
\end{tabular}
\end{table}

The fundamental equation governing Layer-1 gas cost per transaction is:
$$ \text{Cost per tx} = \left(\frac{\text{fixed batch settlement \& proof cost}}{\text{batch size}}\right) + \text{marginal DA cost per tx} $$
Larger batches reduce the amortized fixed cost per transaction but inherently increase queueing latency.
